\documentclass[12pt,a4paper]{article}
\usepackage[margin=2.5cm]{geometry}
\usepackage{booktabs}
\usepackage{longtable}
\usepackage{array}
\usepackage{xcolor}
\usepackage{hyperref}
\usepackage{listings}
\usepackage{amsmath}

\hypersetup{colorlinks=true, linkcolor=blue, urlcolor=blue}

\lstdefinestyle{bnf}{
  basicstyle=\ttfamily\small,
  breaklines=true,
  frame=single,
  backgroundcolor=\color{gray!10},
  columns=flexible,
}

\title{\textbf{Decaf Grammar Specification}\\
       \large CS-471L Compiler Construction Lab\\
       University of Engineering \& Technology, Lahore\\
       Spring 2026}
\author{Compiler Project Team}
\date{}

\begin{document}
\maketitle
\tableofcontents
\newpage

% ------------------------------------------------------------------
\section{Language Overview}
% ------------------------------------------------------------------

Decaf is a strongly-typed, object-oriented language designed for compiler
construction courses. It is a strict subset of Java with a simplified type
system and a fixed set of built-in I/O operations. The language supports:

\begin{itemize}
  \item Primitive types: \texttt{int}, \texttt{double}, \texttt{bool},
        \texttt{string}
  \item One-dimensional arrays of any primitive or class type
  \item Class declarations with optional single inheritance
        (\texttt{extends}) and interface implementation (\texttt{implements})
  \item Interface declarations (method prototype lists only)
  \item Function declarations at global scope and as class members
  \item A full expression language with 9 levels of operator precedence
  \item Control flow: \texttt{if}/\texttt{else}, \texttt{while},
        \texttt{for}, \texttt{return}, \texttt{break}
  \item Built-in operations: \texttt{Print}, \texttt{ReadInteger},
        \texttt{ReadLine}, \texttt{New}, \texttt{NewArray}
\end{itemize}

% ------------------------------------------------------------------
\section{BNF Grammar --- LL(1) Form}
% ------------------------------------------------------------------

The grammar below is the parser-ready, left-recursion-free,
left-factored form used by all three parsers in this project.

\subsection{Program Structure}

\begin{lstlisting}[style=bnf]
Program    ::= DeclList
DeclList   ::= Decl DeclList  |  epsilon
Decl       ::= VarDecl | FuncDecl | ClassDecl | InterfaceDecl
\end{lstlisting}

\subsection{Variable and Function Declarations}

\begin{lstlisting}[style=bnf]
VarDecl    ::= Type id ;

FuncDecl   ::= RetType id ( Formals ) Block

RetType    ::= Type  |  void

Type       ::= int TypeSuffix  |  double TypeSuffix
             | bool TypeSuffix  |  string TypeSuffix
             | id TypeSuffix

TypeSuffix ::= [ ] TypeSuffix  |  epsilon

Formals    ::= ParamList  |  epsilon

ParamList  ::= Param ParamTail
ParamTail  ::= , Param ParamTail  |  epsilon
Param      ::= Type id
\end{lstlisting}

\subsection{Class and Interface Declarations}

\begin{lstlisting}[style=bnf]
ClassDecl      ::= class id ClassBase ClassImpl { FieldList }
ClassBase      ::= extends id  |  epsilon
ClassImpl      ::= implements IdentList  |  epsilon
IdentList      ::= id IdentListTail
IdentListTail  ::= , id IdentListTail  |  epsilon
FieldList      ::= Field FieldList  |  epsilon
Field          ::= VarDecl  |  FuncDecl

InterfaceDecl  ::= interface id { ProtoList }
ProtoList      ::= Proto ProtoList  |  epsilon
Proto          ::= RetType id ( Formals ) ;
\end{lstlisting}

\subsection{Blocks and Statements}

\begin{lstlisting}[style=bnf]
Block      ::= { StmtList }
StmtList   ::= Stmt StmtList  |  epsilon
Stmt       ::= VarDecl | ExprStmt | IfStmt | WhileStmt | ForStmt
             | ReturnStmt | BreakStmt | PrintStmt | Block

ExprStmt   ::= Expr ;
IfStmt     ::= if ( Expr ) Stmt ElsePart
ElsePart   ::= else Stmt  |  epsilon
WhileStmt  ::= while ( Expr ) Stmt
ForStmt    ::= for ( ForInit ; Expr ; ForUpdate ) Stmt
ForInit    ::= Expr  |  epsilon
ForUpdate  ::= Expr  |  epsilon

ReturnStmt ::= return RetExprOpt ;
RetExprOpt ::= Expr  |  epsilon
BreakStmt  ::= break ;
PrintStmt  ::= Print ( ExprList ) ;

ExprList     ::= Expr ExprListTail
ExprListTail ::= , Expr ExprListTail  |  epsilon
\end{lstlisting}

\subsection{Expressions}

\begin{lstlisting}[style=bnf]
Expr       ::= OrExpr AssignTail
AssignTail ::= = Expr  |  epsilon

OrExpr   ::= AndExpr OrTail
OrTail   ::= || AndExpr OrTail  |  epsilon

AndExpr  ::= EqExpr AndTail
AndTail  ::= && EqExpr AndTail  |  epsilon

EqExpr   ::= RelExpr EqTail
EqTail   ::= == RelExpr  |  != RelExpr  |  epsilon

RelExpr  ::= AddExpr RelTail
RelTail  ::= < AddExpr | <= AddExpr | > AddExpr | >= AddExpr | epsilon

AddExpr  ::= MulExpr AddTail
AddTail  ::= + MulExpr AddTail  |  - MulExpr AddTail  |  epsilon

MulExpr  ::= UnaryExpr MulTail
MulTail  ::= * UnaryExpr MulTail | / UnaryExpr MulTail
           | % UnaryExpr MulTail | epsilon

UnaryExpr ::= ! UnaryExpr  |  - UnaryExpr  |  PostExpr

PostExpr  ::= Primary PostTail
PostTail  ::= [ Expr ] PostTail
            | . id CallTail PostTail
            | epsilon

CallTail  ::= ( ActualList )  |  epsilon

Primary   ::= id CallTail | INT_CONST | DOUBLE_CONST | BOOL_CONST
            | STRING_CONST | null | this | ( Expr )
            | ReadInteger ( ) | ReadLine ( )
            | New ( id ) | NewArray ( Expr , Type )

ActualList ::= Expr ExprListTail  |  epsilon
\end{lstlisting}

% ------------------------------------------------------------------
\section{Left-Recursion Removal}
% ------------------------------------------------------------------

\subsection{Array Type Suffix}

The original Decaf specification includes the left-recursive rule:
\[
  \textit{Type} \;\rightarrow\; \textit{int} \;\mid\; \ldots \;\mid\;
  \textit{Type}\;\texttt{[]}\quad\text{(left-recursive)}
\]

This was transformed by introducing \texttt{TypeSuffix}:
\begin{align*}
  \textit{Type}       &\rightarrow \texttt{int}\;\textit{TypeSuffix}
                         \;\mid\; \ldots \;\mid\;
                         \texttt{id}\;\textit{TypeSuffix}\\
  \textit{TypeSuffix} &\rightarrow \texttt{[}\;\texttt{]}\;\textit{TypeSuffix}
                         \;\mid\; \varepsilon
\end{align*}

\subsection{Expression Left Recursion}

The standard left-recursive arithmetic grammar:
\[
  \textit{AddExpr} \;\rightarrow\; \textit{AddExpr}\;\texttt{+}\;\textit{MulExpr}
  \;\mid\; \textit{AddExpr}\;\texttt{-}\;\textit{MulExpr}
  \;\mid\; \textit{MulExpr}
\]
was transformed by introducing right-recursive tail non-terminals:
\begin{align*}
  \textit{AddExpr} &\rightarrow \textit{MulExpr}\;\textit{AddTail}\\
  \textit{AddTail} &\rightarrow \texttt{+}\;\textit{MulExpr}\;\textit{AddTail}
                      \;\mid\; \texttt{-}\;\textit{MulExpr}\;\textit{AddTail}
                      \;\mid\; \varepsilon
\end{align*}

The same pattern applies to all binary operator chains in the expression hierarchy.

% ------------------------------------------------------------------
\section{Operator Precedence Table}
% ------------------------------------------------------------------

\begin{center}
\begin{tabular}{clll}
\toprule
\textbf{Level} & \textbf{Operator(s)} & \textbf{Assoc.} & \textbf{Grammar Rule} \\
\midrule
1 & \texttt{[]}\ \ \texttt{.}               & Left  & PostTail   \\
2 & \texttt{!}\ \ unary \texttt{-}          & Right & UnaryExpr  \\
3 & \texttt{*}\ \ \texttt{/}\ \ \texttt{\%} & Left  & MulTail    \\
4 & \texttt{+}\ \ \texttt{-}                & Left  & AddTail    \\
5 & \texttt{<}\ \ \texttt{<=}\ \ \texttt{>}\ \ \texttt{>=} & Left & RelTail \\
6 & \texttt{==}\ \ \texttt{!=}              & Left  & EqTail     \\
7 & \texttt{\&\&}                           & Left  & AndTail    \\
8 & \texttt{||}                             & Left  & OrTail     \\
9 & \texttt{=}                              & Right & AssignTail \\
\bottomrule
\end{tabular}
\end{center}

% ------------------------------------------------------------------
\section{LL(1) Conflict Analysis}
% ------------------------------------------------------------------

\subsection{Conflict 1 --- Stmt: VarDecl vs ExprStmt on \texttt{id}}

When the lookahead is \texttt{id}, both \textit{VarDecl} (via
\textit{Type} $\rightarrow$ \texttt{id} \textit{TypeSuffix}) and
\textit{ExprStmt} (via \textit{Primary} $\rightarrow$ \texttt{id}
\textit{CallTail}) are candidates. Resolution: a two-token lookahead.
If \texttt{id} is followed by another \texttt{id}, it is a \textit{VarDecl};
otherwise it is an expression statement.

\subsection{Conflict 2 --- Dangling Else}

The standard dangling-else ambiguity is resolved by preferring \textbf{shift}
(binding \texttt{else} to the nearest enclosing \texttt{if}). This is implemented
as a conflict-resolution rule in the SLR(1) table builder and as a greedy match
in the recursive-descent parser.

\subsection{Conflict 3 --- TypeSuffix on \texttt{[}}

State 108 of the SLR(1) automaton encounters both a possible shift (new
\texttt{TypeSuffix}) and a possible reduce. Resolved by reducing, since by
that state the type-declaration context is complete.

\subsection{Conflict 4 --- Shift/Reduce on \texttt{id} in State 1}

Resolved by reducing, which correctly prioritises completing the pending
\textit{DeclList} item before shifting a new \textit{Decl}.

% ------------------------------------------------------------------
\section{Complete Production Index}
% ------------------------------------------------------------------

\begin{longtable}{rp{11cm}}
\toprule
\textbf{\#} & \textbf{Production} \\
\midrule
\endhead
\bottomrule
\endfoot
0  & Program $\rightarrow$ DeclList \\
1  & DeclList $\rightarrow$ Decl DeclList \\
2  & DeclList $\rightarrow$ $\varepsilon$ \\
3  & Decl $\rightarrow$ VarDecl \\
4  & Decl $\rightarrow$ FuncDecl \\
5  & Decl $\rightarrow$ ClassDecl \\
6  & Decl $\rightarrow$ InterfaceDecl \\
7  & VarDecl $\rightarrow$ Type \texttt{id} \texttt{;} \\
8  & FuncDecl $\rightarrow$ RetType \texttt{id} \texttt{(} Formals \texttt{)} Block \\
9  & RetType $\rightarrow$ Type \\
10 & RetType $\rightarrow$ \texttt{void} \\
11 & Type $\rightarrow$ \texttt{int} TypeSuffix \\
12 & Type $\rightarrow$ \texttt{double} TypeSuffix \\
13 & Type $\rightarrow$ \texttt{bool} TypeSuffix \\
14 & Type $\rightarrow$ \texttt{string} TypeSuffix \\
15 & Type $\rightarrow$ \texttt{id} TypeSuffix \\
16 & TypeSuffix $\rightarrow$ \texttt{[} \texttt{]} TypeSuffix \\
17 & TypeSuffix $\rightarrow$ $\varepsilon$ \\
18 & Formals $\rightarrow$ ParamList \\
19 & Formals $\rightarrow$ $\varepsilon$ \\
20 & ParamList $\rightarrow$ Param ParamTail \\
21 & ParamTail $\rightarrow$ \texttt{,} Param ParamTail \\
22 & ParamTail $\rightarrow$ $\varepsilon$ \\
23 & Param $\rightarrow$ Type \texttt{id} \\
24 & ClassDecl $\rightarrow$ \texttt{class} \texttt{id} ClassBase ClassImpl \texttt{\{} FieldList \texttt{\}} \\
25 & ClassBase $\rightarrow$ \texttt{extends} \texttt{id} \\
26 & ClassBase $\rightarrow$ $\varepsilon$ \\
27 & ClassImpl $\rightarrow$ \texttt{implements} IdentList \\
28 & ClassImpl $\rightarrow$ $\varepsilon$ \\
29 & IdentList $\rightarrow$ \texttt{id} IdentListTail \\
30 & IdentListTail $\rightarrow$ \texttt{,} \texttt{id} IdentListTail \\
31 & IdentListTail $\rightarrow$ $\varepsilon$ \\
32 & FieldList $\rightarrow$ Field FieldList \\
33 & FieldList $\rightarrow$ $\varepsilon$ \\
34 & Field $\rightarrow$ VarDecl \\
35 & Field $\rightarrow$ FuncDecl \\
36 & InterfaceDecl $\rightarrow$ \texttt{interface} \texttt{id} \texttt{\{} ProtoList \texttt{\}} \\
37 & ProtoList $\rightarrow$ Proto ProtoList \\
38 & ProtoList $\rightarrow$ $\varepsilon$ \\
39 & Proto $\rightarrow$ RetType \texttt{id} \texttt{(} Formals \texttt{)} \texttt{;} \\
40 & Block $\rightarrow$ \texttt{\{} StmtList \texttt{\}} \\
41 & StmtList $\rightarrow$ Stmt StmtList \\
42 & StmtList $\rightarrow$ $\varepsilon$ \\
43--51 & Stmt $\rightarrow$ VarDecl $\mid$ ExprStmt $\mid$ IfStmt $\mid$ WhileStmt $\mid$ ForStmt $\mid$ ReturnStmt $\mid$ BreakStmt $\mid$ PrintStmt $\mid$ Block \\
52 & ExprStmt $\rightarrow$ Expr \texttt{;} \\
53 & IfStmt $\rightarrow$ \texttt{if} \texttt{(} Expr \texttt{)} Stmt ElsePart \\
54 & ElsePart $\rightarrow$ \texttt{else} Stmt \\
55 & ElsePart $\rightarrow$ $\varepsilon$ \\
56 & WhileStmt $\rightarrow$ \texttt{while} \texttt{(} Expr \texttt{)} Stmt \\
57 & ForStmt $\rightarrow$ \texttt{for} \texttt{(} ForInit \texttt{;} Expr \texttt{;} ForUpdate \texttt{)} Stmt \\
58 & ForInit $\rightarrow$ Expr \\
59 & ForInit $\rightarrow$ $\varepsilon$ \\
60 & ForUpdate $\rightarrow$ Expr \\
61 & ForUpdate $\rightarrow$ $\varepsilon$ \\
62 & ReturnStmt $\rightarrow$ \texttt{return} RetExprOpt \texttt{;} \\
63 & RetExprOpt $\rightarrow$ Expr \\
64 & RetExprOpt $\rightarrow$ $\varepsilon$ \\
65 & BreakStmt $\rightarrow$ \texttt{break} \texttt{;} \\
66 & PrintStmt $\rightarrow$ \texttt{Print} \texttt{(} ExprList \texttt{)} \texttt{;} \\
67 & ExprList $\rightarrow$ Expr ExprListTail \\
68 & ExprListTail $\rightarrow$ \texttt{,} Expr ExprListTail \\
69 & ExprListTail $\rightarrow$ $\varepsilon$ \\
70 & Expr $\rightarrow$ OrExpr AssignTail \\
71 & AssignTail $\rightarrow$ \texttt{=} Expr \\
72 & AssignTail $\rightarrow$ $\varepsilon$ \\
73 & OrExpr $\rightarrow$ AndExpr OrTail \\
74 & OrTail $\rightarrow$ \texttt{||} AndExpr OrTail \\
75 & OrTail $\rightarrow$ $\varepsilon$ \\
76 & AndExpr $\rightarrow$ EqExpr AndTail \\
77 & AndTail $\rightarrow$ \texttt{\&\&} EqExpr AndTail \\
78 & AndTail $\rightarrow$ $\varepsilon$ \\
79 & EqExpr $\rightarrow$ RelExpr EqTail \\
80 & EqTail $\rightarrow$ \texttt{==} RelExpr \\
81 & EqTail $\rightarrow$ \texttt{!=} RelExpr \\
82 & EqTail $\rightarrow$ $\varepsilon$ \\
83 & RelExpr $\rightarrow$ AddExpr RelTail \\
84--88 & RelTail $\rightarrow$ \texttt{<} AddExpr $\mid$ \texttt{<=} AddExpr $\mid$ \texttt{>} AddExpr $\mid$ \texttt{>=} AddExpr $\mid$ $\varepsilon$ \\
89 & AddExpr $\rightarrow$ MulExpr AddTail \\
90 & AddTail $\rightarrow$ \texttt{+} MulExpr AddTail \\
91 & AddTail $\rightarrow$ \texttt{-} MulExpr AddTail \\
92 & AddTail $\rightarrow$ $\varepsilon$ \\
93 & MulExpr $\rightarrow$ UnaryExpr MulTail \\
94 & MulTail $\rightarrow$ \texttt{*} UnaryExpr MulTail \\
95 & MulTail $\rightarrow$ \texttt{/} UnaryExpr MulTail \\
96 & MulTail $\rightarrow$ \texttt{\%} UnaryExpr MulTail \\
97 & MulTail $\rightarrow$ $\varepsilon$ \\
98 & UnaryExpr $\rightarrow$ \texttt{!} UnaryExpr \\
99 & UnaryExpr $\rightarrow$ \texttt{-} UnaryExpr \\
100 & UnaryExpr $\rightarrow$ PostExpr \\
101 & PostExpr $\rightarrow$ Primary PostTail \\
102 & PostTail $\rightarrow$ \texttt{[} Expr \texttt{]} PostTail \\
103 & PostTail $\rightarrow$ \texttt{.} \texttt{id} CallTail PostTail \\
104 & PostTail $\rightarrow$ $\varepsilon$ \\
105 & CallTail $\rightarrow$ \texttt{(} ActualList \texttt{)} \\
106 & CallTail $\rightarrow$ $\varepsilon$ \\
107--118 & Primary productions (see text) \\
119 & ActualList $\rightarrow$ Expr ExprListTail \\
120 & ActualList $\rightarrow$ $\varepsilon$ \\
\end{longtable}

\end{document}
